\documentclass[11pt, a4paper]{article}

% Gemeinsame Style-Definitionen (TEXINPUTS muss templates/ enthalten — siehe scripts/build_tex.py)
\input{bfw-style.tex}

\title{\vspace*{-1.5cm}\Huge\bfseries\color{bfwPrimaryDark}Session~7: Präsentation vorbereiten \& planen\\[0.4em]
  \large\color{bfwInkSoft}\textnormal{Teamarbeit, Planung und Medienwahl für die Betriebspräsentation bei JIKU IT-Solutions}}
\author{\textsf{IT Lernfeld 1 --- Das Unternehmen und die eigene Rolle im Betrieb beschreiben}}
\date{}

\begin{document}
\maketitle
\thispagestyle{empty}

\vspace{1em}
\begin{merksatz}[title={\bfseries Worum geht es in dieser Session?}]
Die Aufgabe lautet: Den Ausbildungsbetrieb \emph{JIKU IT-Solutions GmbH} einem externen Publikum
präsentieren. Dafür braucht es mehr als gute Folien --- es braucht strukturierte Teamarbeit, eine
klare Zielgruppenanalyse, einen roten Faden im Aufbau und eine durchdachte Medienwahl.
In dieser Session erarbeiten wir alle Planungsschritte, die eine professionelle Präsentation erst
ermöglichen.
\end{merksatz}

\tableofcontents
\vspace{1em}
\hrule
\vspace{1em}

\section{Die Präsentation in Teamarbeit vorbereiten (Kapitel 1.5.1)}

\subsection{Gruppenarbeit und Projektarbeit}

Bei der \textbf{Gruppenarbeit} arbeiten alle Personen gemeinsam und weitgehend gleichwertig an
einer Aufgabe. Bei der \textbf{Projektarbeit} sind Personen mit \emph{unterschiedlichen Kompetenzen}
eingebunden; eine Projektleitung koordiniert die Zusammenarbeit und sorgt dafür, dass das gemeinsame
Ziel mit hoher Interaktion erreicht wird.

\begin{merksatz}[title={\bfseries Merke: Projektteam}]
Ein \textbf{Projektteam} ist eine zeitlich begrenzte Arbeitsgruppe, in der Personen mit
unterschiedlichen Fähigkeiten gemeinsam an einem klar definierten Sonderauftrag arbeiten.
Neben fachlicher Kompetenz trägt ein ausgeprägtes Wir-Gefühl wesentlich zum Gelingen bei.
\end{merksatz}

\subsection{Fünf Erfolgsfaktoren für ein gutes Team}

\begin{enumerate}
  \item \textbf{Klares Ziel:} Alle müssen die gleiche Richtung kennen. Zielformulierungen brauchen
        daher besondere Aufmerksamkeit --- verständlich und konkret.
  \item \textbf{Richtige Teamgröße:} Fünf bis acht Personen gelten als optimal für kreative und
        intensive Arbeit. Bei größeren Projekten werden Teilprojekte gebildet.
  \item \textbf{Unterschiedliche Teampersönlichkeiten:} Homogene Teams erzeugen weniger Reibung, aber
        auch weniger neue Ideen. Nützliche Rollen: Analytiker, Querdenker, Kreative, Innovatoren,
        Pragmatiker, Macher, Teamer, Organisatoren.
  \item \textbf{Akzeptierte Projektleitung:} Die Projektleitung muss Verantwortung übernehmen, das
        große Ganze im Blick behalten, moderieren und besonnen agieren.
  \item \textbf{Funktionierende Kommunikation:} Alle Projektbeteiligten müssen gut informiert sein und
        zeitnah die richtigen Maßnahmen abstimmen können.
\end{enumerate}

\subsection{Teamkonflikte erkennen und steuern}

Zwei verbreitete Gruppenphänomene gefährden den Projekterfolg:

\begin{center}
\begin{tabular}{p{4.5cm} p{8cm}}
\toprule
\textbf{Phänomen} & \textbf{Beschreibung und Folge}\\
\midrule
Soziales Faulenzen & Einzelne reduzieren ihre Leistung, weil sie sich in der Gruppe „verstecken"
                     können. Wenn es ohne Konsequenzen bleibt, sinkt die Leistung der gesamten Gruppe.\\
\addlinespace
Trittbrettfahren & Jemand profitiert vom Gruppenergebnis, ohne einen fairen Beitrag zu leisten.
                   Leistungsträger ziehen sich zurück, der Projekterfolg gerät in Gefahr.\\
\bottomrule
\end{tabular}
\end{center}

Gegenmaßnahmen: Kooperationsbereitschaft fördern, feste Regeln vereinbaren, Aufgaben mit
namentlicher Verantwortung verteilen, Zwischenberichte einplanen.

\subsection{Phasenmodell der Teamentwicklung nach Bruce Tuckman}

Der US-amerikanische Psychologe Bruce Tuckman beschrieb 1965 folgende Entwicklungsphasen für Teams:

\begin{center}
\begin{tabular}{p{2.5cm} p{3.5cm} p{7cm}}
\toprule
\textbf{Phase} & \textbf{Name} & \textbf{Merkmale}\\
\midrule
Phase 1 & Forming (Einstieg) & Kennenlernen, höflicher Umgang, erste Regeln werden vereinbart\\
\addlinespace
Phase 2 & Storming (Streit) & Unterschwellige Konflikte, Cliquenbildung, Aufgaben treten in den Hintergrund\\
\addlinespace
Phase 3 & Norming (Regelung) & Regeln werden diskutiert, gefunden und anerkannt; Rollen festigen sich\\
\addlinespace
Phase 4 & Performing (Leistung) & Zielorientiert, flexibel, ideen­reich; gegenseitige Wertschätzung\\
\addlinespace
Phase 5 & Adjourning (Auflösung) & Abschluss des Projekts, Auflösung des Teams\\
\bottomrule
\end{tabular}
\end{center}

\begin{merksatz}[title={\bfseries Merke: Teamphasen bewusst begleiten}]
Projektleitungen sollten damit rechnen, dass Teams alle Phasen durchlaufen. Durch Vorgespräche,
klare Regeln und gezielte Feedbacksitzungen lässt sich die Entwicklung zu einer leistungsfähigen
Gruppe aktiv unterstützen.
\end{merksatz}

\subsection{Die IPERKA-Methode}

IPERKA steht für die sechs Phasen der vollständigen Handlung und strukturiert die Arbeit eines
Projektteams:

\begin{center}
\begin{tabular}{p{1.5cm} p{3cm} p{8.5cm}}
\toprule
\textbf{Buchst.} & \textbf{Phase} & \textbf{Inhalt}\\
\midrule
I & Informieren & Aufgabe verstehen; Informationen recherchieren oder erfragen\\
\addlinespace
P & Planen & Arbeitsschritte, Zeitbedarf, Hilfsmittel, Teilziele und Termine festlegen\\
\addlinespace
E & Entscheiden & Welcher Plan wird umgesetzt? Wer übernimmt die Leitung?\\
\addlinespace
R & Realisieren & Den Arbeits- und Maßnahmenplan wie vereinbart umsetzen\\
\addlinespace
K & Kontrollieren & Ergebnisse prüfen — Selbst- und Fremdkontrolle anhand eines Kriterienbogens\\
\addlinespace
A & Auswerten & Mit dem Auftraggeber reflektieren: Was lief gut? Was ist zu verbessern?\\
\bottomrule
\end{tabular}
\end{center}

\section{Die Präsentation planen (Kapitel 1.5.2)}

\subsection{Präsentationsbewertung und Anforderungen}

Bevor Inhalte ausgewählt werden, lohnt es sich, die \emph{Bewertungskriterien} zu kennen, die an
die Präsentation angelegt werden. Für die mündliche Abschlussprüfung IT-Berufe gilt beispielsweise
folgender Kriterienkatalog der IHK:

\begin{center}
\begin{tabular}{p{5.5cm} p{7cm}}
\toprule
\textbf{Bewertungsbereich} & \textbf{Inhalt (max. je 10 Punkte)}\\
\midrule
Aufbau und inhaltliche Struktur & Gliederung, Logik, Zielorientierung\\
\addlinespace
Sprachliche Gestaltung & Ausdrucksweise, Satzbau, Stil\\
\addlinespace
Zielgruppengerechte Darstellung & Medieneinsatz, Visualisierung, Körpersprache\\
\midrule
\textbf{Gesamt} & \textbf{max. 30 Punkte}\\
\bottomrule
\end{tabular}
\end{center}

\subsection{Rahmenbedingungen klären}

Vor der eigentlichen Inhaltserstellung müssen die organisatorischen Rahmenbedingungen bekannt sein:

\begin{center}
\begin{tabular}{p{3.5cm} p{9.5cm}}
\toprule
\textbf{Faktor} & \textbf{Leitfragen}\\
\midrule
Kosten & Wie hoch ist das Budget für die Präsentation (Druck, Technik)?\\
\addlinespace
Zeit & Wie viel Zeit steht für Vortrag, Aufbau, Einleitung und Schlussgespräch zur Verfügung? Tageszeit?\\
\addlinespace
Vorkenntnisse & Über welches Vorwissen verfügen die Teilnehmer? Gibt es Überschneidungen mit anderen Vorträgen?\\
\addlinespace
Interessen & Eher inhaltliche Tiefe oder Breite? Wie viel Zeit für Diskussion und Fragen?\\
\addlinespace
Raum & Ausstattung mit Strom, Medien, Akustik und Sichtverhältnissen für alle Teilnehmer?\\
\addlinespace
Unterlagen & Welche Unterlagen sind bereit­zustellen (Handout, Dateien, Feedback­bogen, Protokoll)?\\
\bottomrule
\end{tabular}
\end{center}

\subsection{Informationsbeschaffung und Stoffauswahl}

Zur Informationssammlung stehen verschiedene Methoden zur Verfügung: Brainstorming, Vor-Ort-Recherche,
Internetrecherche, Fachbuchrecherche, Befragungen, Bibliotheken, Statistiken und eigene Erfahrungen.

\textbf{Leitfragen für die Stoffauswahl:}

\begin{itemize}
  \item Passt der Inhalt zum Thema?
  \item Ist er wichtig für das Ziel der Präsentation?
  \item Interessiert er die Teilnehmer?
  \item Wurden geeignete Beispiele gefunden?
  \item Gibt es Zeitprobleme?
  \item Können Fakten reduziert oder zusammengefasst werden?
\end{itemize}

\begin{merksatz}[title={\bfseries KISS-Regel und 10-20-30-Regel}]
\textbf{KISS:} Keep It Short and Simple --- halte es kurz und einfach! Inhalte auf das Wesentliche
verdichten.

\textbf{10-20-30-Regel} (bekannte Faustregel, nach Guy Kawasaki): Als Orientierung gilt ---
maximal \textbf{10 Folien}, Vortragsdauer höchstens \textbf{20 Minuten} und eine Mindestschriftgröße
von \textbf{30 Punkt} für gut lesbare Folien.
\end{merksatz}

\subsection{Aufbau: Einstieg --- Hauptteil --- Schluss}

Eine gut gegliederte Präsentation folgt dem klassischen dreiteiligen Aufbau:

\begin{center}
\begin{tabular}{p{3cm} p{5.5cm} p{4cm}}
\toprule
\textbf{Abschnitt} & \textbf{Inhalt und Funktion} & \textbf{Zeitanteil\textsuperscript{*}}\\
\midrule
Einstieg & Interesse wecken, Thema nennen, Überblick geben — die ersten Sätze prägen sich ein & ca. 10--15\,\%\\
\addlinespace
Hauptteil & Sachlogisch gegliederte Kernaussagen; Spannungsbogen halten; Beispiele einsetzen & ca. 70--75\,\%\\
\addlinespace
Schluss & Zusammenfassung, Fazit, Ausblick — die letzten Sätze bleiben in Erinnerung & ca. 10--15\,\%\\
\bottomrule
\end{tabular}

\smallskip
\footnotesize\textsuperscript{*}Praxisübliche Richtwerte als Faustregel; nicht aus dem Lehrwerk.
\end{center}

Nützliche Ordnungsprinzipien für den Aufbau des Hauptteils:
\begin{itemize}
  \item Vom Problem zur Problemlösung
  \item Vom Allgemeinen zum Besonderen
  \item Vom Leichten zum Schweren
  \item Vom Bekannten zum Unbekannten
\end{itemize}

\subsection{Medienwahl und Visualisierung}

\begin{center}
\begin{tabular}{p{3.5cm} p{9.5cm}}
\toprule
\textbf{Kategorie} & \textbf{Beispiele}\\
\midrule
Flächen/Träger & Flip-Chart, Stell-/Pinnwand, Plakate, Tafel, Whiteboard, Beamer-Präsentation\\
\addlinespace
Objekte & Karten, Zeichnungen, Bilder, Diagramme, Videos, Symbole\\
\addlinespace
Präsentationsutensilien & Stifte, Marker, Pointer, Moderatorenkoffer\\
\addlinespace
Vortragsextras & Handout, Skript, Rollenspiel, Rundgang, Vorführung\\
\addlinespace
Multimedia-Tools & Microsoft PowerPoint, Google Slides, Prezi, Apple Keynote, OpenOffice Impress\\
\bottomrule
\end{tabular}
\end{center}

\begin{merksatz}[title={\bfseries Merke: Behaltensquote und Sinnesvielfalt}]
Je mehr Sinne angesprochen werden, desto höher die Behaltensquote:
\textbf{nur hören: 20\,\%} --- \textbf{nur sehen: 30\,\%} --- \textbf{hören + sehen: 50\,\%} ---
\textbf{hören + sehen + darüber sprechen: 70\,\%} ---
\textbf{hören + sehen + sprechen + selber tun: 90\,\%}.
Gute Präsentationen sprechen möglichst viele Sinne gleichzeitig an.
\end{merksatz}

\subsection{Präsentationsregeln: Sprache und Auftreten}

\begin{center}
\begin{tabular}{p{3.5cm} p{9.5cm}}
\toprule
\textbf{Bereich} & \textbf{Regeln}\\
\midrule
Sprache & Frei, anschaulich und lebendig; natürlich und deutlich; in angemessener Lautstärke;
          normale Sprechgeschwindigkeit; Tonlage variieren, um Wichtiges zu betonen\\
\addlinespace
Körpersprache & Aufrechte Haltung; natürliche Gestik und Mimik; Blickkontakt zum Publikum aufnehmen;
                Zuhörer einbeziehen\\
\addlinespace
Nervosität & Stolpersteine im Vorfeld beseitigen; positive innere Haltung; Atemübung: einatmen,
             4 Sekunden halten, langsam ausatmen\\
\addlinespace
Einstieg & Laut und deutlich beginnen; schnell Blickkontakt aufnehmen; interessanten ersten Satz
           formulieren\\
\addlinespace
Timing & Angekündigte Vortragsdauer nicht überschreiten; kurze Sätze und Sprechpausen einplanen;
         Generalprobe mit Zeitnahme durchführen\\
\bottomrule
\end{tabular}
\end{center}

\begin{merksatz}[title={\bfseries Zusammenfassung: Kernbegriffe dieser Session}]
\textbf{IPERKA:} Informieren, Planen, Entscheiden, Realisieren, Kontrollieren, Auswerten --- die
vollständige Handlung im Projektteam.~·~
\textbf{Teamentwicklung nach Tuckman:} Forming → Storming → Norming → Performing → Adjourning.~·~
\textbf{KISS-Regel:} Keep It Short and Simple.~·~
\textbf{10-20-30-Regel:} max.\ 10 Folien · max.\ 20 Minuten · min.\ 30 Punkt Schriftgröße.~·~
\textbf{Aufbau:} Einstieg --- Hauptteil --- Schluss mit rotem Faden.~·~
\textbf{Behaltensquote:} steigt von 20\,\% (nur hören) auf 90\,\% (alle Sinne).
\end{merksatz}

\end{document}
